% UPMurphi Algorithm Flow Overview
% Compile: xelatex algorithm_flowchart.tex

\documentclass[border=12pt,tikz]{standalone}

\usepackage[UTF8]{ctex}
\usepackage{tikz}
\usepackage{amsmath}
\usetikzlibrary{
  shapes.geometric,
  arrows.meta,
  positioning,
  calc,
  fit,
  backgrounds,
  decorations.pathreplacing
}

% ── colors ──
\definecolor{compilebg}{HTML}{E8F0FE}
\definecolor{compileborder}{HTML}{4285F4}
\definecolor{phase1bg}{HTML}{FFF3E0}
\definecolor{phase1border}{HTML}{FB8C00}
\definecolor{phase2bg}{HTML}{E8F5E9}
\definecolor{phase2border}{HTML}{43A047}
\definecolor{phase3bg}{HTML}{F3E5F5}
\definecolor{phase3border}{HTML}{8E24AA}
\definecolor{phase4bg}{HTML}{E0F7FA}
\definecolor{phase4border}{HTML}{00ACC1}
\definecolor{phase56bg}{HTML}{FCE4EC}
\definecolor{phase56border}{HTML}{E53935}
\definecolor{algbfs}{HTML}{1565C0}
\definecolor{algdfs}{HTML}{2E7D32}
\definecolor{algastar}{HTML}{C62828}

% ── helper command for phase title ──
\newcommand{\phasetitle}[2]{{\small\bfseries\color{#1}#2}}

\begin{document}
\begin{tikzpicture}[
  % ── node styles ──
  iobox/.style={
    rectangle, rounded corners=3pt,
    draw=compileborder, line width=1pt, fill=compileborder!8,
    text width=20em, align=center,
    inner sep=6pt, font=\small
  },
  phasebox/.style 2 args={
    rectangle, rounded corners=5pt,
    draw=#1, line width=1.2pt, fill=#2,
    text width=26em, align=left,
    inner sep=10pt, font=\small
  },
  algobox/.style={
    rectangle, rounded corners=2pt,
    draw=#1, line width=0.8pt, fill=#1!12,
    minimum height=2.2em, minimum width=5.2em,
    align=center, font=\small\bfseries, text=#1
  },
  myarrow/.style={
    -{Stealth[length=6pt,width=5pt]},
    line width=1pt, color=black!70
  },
  complabel/.style={
    font=\scriptsize\itshape, text=black!70
  },
]

% ═══════════════════════════════════════════════════
%   Compilation chain (top)
% ═══════════════════════════════════════════════════

\node[iobox] (pddl) {
  \textbf{PDDL+ Domain} + \textbf{Problem}
};

\node[iobox, below=0.55cm of pddl] (model) {
  UPMurphi Model \texttt{(.m)}
};

\node[iobox, below=0.55cm of model] (cppsrc) {
  C++ Source \texttt{(.cpp)}
};

\node[iobox, below=0.55cm of cppsrc] (planner) {
  \textbf{Executable Planner}
};

\draw[myarrow] (pddl)  -- node[right, complabel] {\texttt{pddl2upm}} (model);
\draw[myarrow] (model)  -- node[right, complabel] {\texttt{upmc}}     (cppsrc);
\draw[myarrow] (cppsrc) -- node[right, complabel] {\texttt{g++ -m32}} (planner);

% dashed compile group
\begin{scope}[on background layer]
  \node[fit=(pddl)(planner),
        draw=compileborder!50, dashed, line width=0.8pt,
        rounded corners=8pt, fill=compilebg,
        inner sep=12pt] (compilegroup) {};
\end{scope}
\node[above=2pt of compilegroup.north,
      font=\footnotesize\bfseries, text=compileborder]
  {编译工具链 (Compilation Toolchain)};

% ═══════════════════════════════════════════════════
%   Phase 1: State Space Exploration
% ═══════════════════════════════════════════════════

\node[phasebox={phase1border}{phase1bg}, below=1.4cm of planner] (phase1) {
  \phasetitle{phase1border}{Phase 1: 状态空间探索 (State Space Exploration)}%
  \\[24pt]% leave room for the algorithm boxes drawn on top
  对每个状态调用 \textsc{AllNextStates}():\\[2pt]
  \quad$\triangleright$ 施加所有 enabled rules\\
  \quad$\triangleright$ \textsc{Add}(nextstate, transition)\\
  \quad$\triangleright$ 记录转移 (from $\rightarrow$ to, rule)\\[4pt]
  \textbf{输出:} 状态集合, 转移文件, 目标文件
};

% three algorithm boxes inside phase 1
\node[algobox=algbfs,   anchor=south] at ($(phase1.north)+(-7em,-1.6em)$)  (bfs)   {BFS\\\scriptsize FIFO Queue};
\node[algobox=algdfs,   anchor=south] at ($(phase1.north)+(0em,-1.6em)$)   (dfs)   {DFS\\\scriptsize Stack};
\node[algobox=algastar, anchor=south] at ($(phase1.north)+(7em,-1.6em)$)   (astar) {A*\\\scriptsize MinHeap\\[-2pt]\scriptsize $f\!=\!g\!+\!h$};

% merge lines from three algo boxes
\coordinate (algmerge) at ($(phase1.north)+(0em,-4em)$);
\draw[line width=0.6pt, color=black!50] (bfs.south)   |- (algmerge);
\draw[line width=0.6pt, color=black!50] (dfs.south)   --  (algmerge);
\draw[line width=0.6pt, color=black!50] (astar.south) |- (algmerge);
\filldraw[black!50] (algmerge) circle (1.5pt);

% arrow from planner to phase 1
\draw[myarrow] (planner.south) -- (phase1.north);

% ═══════════════════════════════════════════════════
%   Phase 2: Build Dynamics
% ═══════════════════════════════════════════════════

\node[phasebox={phase2border}{phase2bg}, below=0.9cm of phase1] (phase2) {
  \phasetitle{phase2border}{Phase 2: 构建反转图 (Build Inverted Graph)}%
  \\[4pt]
  将 from$\,\rightarrow\,$to 转移构建为 to$\,\rightarrow\,$from 反转邻接表\\
  用于 Phase~3 的反向路径搜索
};

\draw[myarrow] (phase1.south) -- (phase2.north);

% ═══════════════════════════════════════════════════
%   Phase 3: Find Paths
% ═══════════════════════════════════════════════════

\node[phasebox={phase3border}{phase3bg}, below=0.9cm of phase2] (phase3) {
  \phasetitle{phase3border}{Phase 3: 反向路径搜索 (Backward Dijkstra)}%
  \\[4pt]
  从 Goal 状态出发, 反向 Dijkstra 到所有可达状态\\
  为每个可控状态选择最优动作\\[4pt]
  搜索策略:\\
  \quad$\bullet$ \textbf{Feasible:} 任意一条路径\\
  \quad$\bullet$ \textbf{Optimal:} 最短/最轻路径\\
  \quad$\bullet$ \textbf{Universal:} 所有状态的通用策略
};

\draw[myarrow] (phase2.south) -- (phase3.north);

% ═══════════════════════════════════════════════════
%   Phase 4: Collect Plans
% ═══════════════════════════════════════════════════

\node[phasebox={phase4border}{phase4bg}, below=0.9cm of phase3] (phase4) {
  \phasetitle{phase4border}{Phase 4: 收集计划 (Collect Plans)}%
  \\[4pt]
  从每个初始状态沿最优动作链追踪到目标\\
  生成完整计划序列
};

\draw[myarrow] (phase3.south) -- (phase4.north);

% ═══════════════════════════════════════════════════
%   Phase 5 & 6: Output + Validate
% ═══════════════════════════════════════════════════

\node[phasebox={phase56border}{phase56bg}, below=0.9cm of phase4] (phase56) {
  \phasetitle{phase56border}{Phase 5: 输出结果 (Output Results)}%
  \\[4pt]
  格式: PDDL+ / Text / CSV / Raw\\[6pt]
  \phasetitle{phase56border}{Phase 6: 自动验证 (Validate Results) \normalfont\itshape\small [可选]}%
  \\[4pt]
  调用 VAL 验证器验证生成的计划
};

\draw[myarrow] (phase4.south) -- (phase56.north);

% ═══════════════════════════════════════════════════
%   Right-side brace annotation
% ═══════════════════════════════════════════════════

\coordinate (braceTop) at ($(phase1.north east)+(0.8em,0)$);
\coordinate (braceBot) at ($(phase56.south east)+(0.8em,0)$);

\draw[decorate, decoration={brace, amplitude=8pt, mirror},
      line width=0.8pt, color=black!40]
  (braceTop) -- (braceBot)
  node[midway, right=10pt, align=left, font=\footnotesize, text=black!55] {
    运行时\\(Runtime)
  };

\end{tikzpicture}
\end{document}
